% Памятка: программирование в симуляторе Пионера.
% Сборка: latexmk -xelatex pioneer-sim-guide.tex
\documentclass[11pt,a4paper]{article}

\usepackage{fontspec}
\usepackage{polyglossia}
\setdefaultlanguage{russian}
\setotherlanguage{english}
\setmainfont{Noto Sans}
\setmonofont{Noto Sans Mono}[Scale=0.9]
\newfontfamily\cyrillicfont{Noto Sans}
\newfontfamily\cyrillicfonttt{Noto Sans Mono}[Scale=0.9]

\usepackage[margin=18mm,top=16mm,bottom=18mm]{geometry}
\usepackage{graphicx}
\usepackage{xcolor}
\usepackage{tikz}
\usetikzlibrary{babel}
\usepackage{listings}
\usepackage[most]{tcolorbox}
\usepackage{booktabs}
\usepackage{tabularx}
\usepackage{enumitem}
\usepackage{titlesec}
\usepackage{fancyhdr}
\usepackage[hidelinks]{hyperref}
\usepackage{multicol}
\usepackage{caption}

\definecolor{accent}{HTML}{FF6B00}
\definecolor{ink}{HTML}{1F2328}
\definecolor{muted}{HTML}{5F6670}
\definecolor{codebg}{HTML}{F6F7F9}
\definecolor{codeframe}{HTML}{E3E6EA}
\definecolor{kw}{HTML}{7C3AED}
\definecolor{fn}{HTML}{E85D04}
\definecolor{str}{HTML}{0F766E}
\definecolor{cmt}{HTML}{6B7280}
\definecolor{tipbg}{HTML}{EEF6FF}
\definecolor{tipframe}{HTML}{2563EB}
\definecolor{warnbg}{HTML}{FFF4EC}

\color{ink}
\setlength{\parindent}{0pt}
\setlength{\emergencystretch}{3em}
\setlength{\parskip}{5pt}
\captionsetup{font={small,color=muted},labelformat=empty,skip=4pt}

\titleformat{\section}{\Large\bfseries\color{ink}}{\color{accent}\thesection}{0.6em}{}
\titleformat{\subsection}{\large\bfseries\color{ink}}{}{0em}{}
\titlespacing*{\section}{0pt}{14pt}{6pt}
\titlespacing*{\subsection}{0pt}{10pt}{4pt}

\pagestyle{fancy}
\fancyhf{}
\renewcommand{\headrulewidth}{0pt}
\fancyfoot[L]{\small\color{muted}Симулятор Пионера · simulator.sakoryagin.ru}
\fancyfoot[R]{\small\color{muted}\thepage}

\lstdefinestyle{base}{
    basicstyle=\ttfamily\small,
    backgroundcolor=\color{codebg},
    frame=single, rulecolor=\color{codeframe}, framesep=6pt, xleftmargin=6pt, xrightmargin=6pt,
    keywordstyle=\color{kw}\bfseries,
    commentstyle=\color{cmt},
    stringstyle=\color{str},
    showstringspaces=false, columns=fullflexible, keepspaces=true,
    breaklines=true, aboveskip=6pt, belowskip=6pt,
    extendedchars=true
}
\lstdefinelanguage{lua}{
    morekeywords={function,end,if,then,else,elseif,for,do,while,local,return,and,or,not,nil,true,false,repeat,until,in},
    morecomment=[l]{--}, morestring=[b]", sensitive=true
}
\lstdefinestyle{lua}{style=base, language=lua,
    emph={ap,Ev,Timer,Sensors,Ledbar,camera,callback,sleep,time}, emphstyle=\color{fn}}
\lstdefinestyle{py}{style=base, language=Python,
    morekeywords={True,False,None,pass},
    emph={pioneer,Pioneer,Camera,time}, emphstyle=\color{fn}}

\newtcolorbox{tip}[1][Совет]{enhanced, colback=tipbg, colframe=tipframe, boxrule=0pt, leftrule=3pt,
    arc=2pt, fonttitle=\bfseries\small, coltitle=tipframe, colbacktitle=tipbg, title={#1},
    attach title to upper={\ \ }, fontupper=\small, left=8pt, top=4pt, bottom=4pt}
\newtcolorbox{warn}[1][Осторожно]{enhanced, colback=warnbg, colframe=accent, boxrule=0pt, leftrule=3pt,
    arc=2pt, fonttitle=\bfseries\small, coltitle=accent, colbacktitle=warnbg, title={#1},
    attach title to upper={\ \ }, fontupper=\small, left=8pt, top=4pt, bottom=4pt}

\newcommand{\shot}[2][\linewidth]{\includegraphics[width=#1]{img/#2}}
% Клавиша — без TikZ: TikZ внутри подписей и списков ломается о русские активные символы.
\newcommand{\key}[1]{{\setlength{\fboxsep}{1.5pt}\fcolorbox{codeframe}{codebg}{\small\ttfamily #1}}}
\newcommand{\ui}[1]{«\textbf{#1}»}
% Номерок-выноска на скриншоте: \callout{x}{y}{n}, x и y — доли ширины и высоты.
\newcommand{\callout}[3]{\node[circle,fill=accent,text=white,font=\small\bfseries,inner sep=1.6pt,minimum size=15pt,draw=white,line width=1pt] at (#1,#2) {#3};}
% Номер в списке под скриншотом — того же цвета, что выноски на нём.
\newcommand{\bubble}[1]{{\setlength{\fboxsep}{1.6pt}\colorbox{accent}{\makebox[0.9em]{\color{white}\scriptsize\bfseries #1}}}}

\begin{document}

% ---------------------------------------------------------------- титул
\begin{tcolorbox}[enhanced, colback=ink, colframe=ink, arc=4pt, left=14pt, right=14pt, top=12pt, bottom=12pt]
    {\color{accent}\small\bfseries ПАМЯТКА}\\[2pt]
    {\color{white}\huge\bfseries\raggedright\hyphenpenalty=10000 Программирование в~симуляторе Пионера\par}\vspace{6pt}
    {\color{white!80!black}Lua, Python и Blockly на одной странице: как запустить первую программу за пять минут
    и не наступить на самые частые грабли.\\[4pt]
    Адрес: \textbf{\color{white}https://simulator.sakoryagin.ru}}
\end{tcolorbox}

\begin{multicols}{2}
\small
\textbf{Что внутри}
\begin{enumerate}[leftmargin=*,itemsep=0pt]
    \item Быстрый старт
    \item Первая программа на Lua
    \item Координаты: где точка
    \item Несколько точек подряд
    \item Таймеры, светодиоды, датчики
    \item Python в браузере
    \item Blockly: блоки вместо кода
    \item Камеры дрона
    \item Python со своего компьютера
    \item Частые ошибки
    \item Шпаргалка: Lua и Python
\end{enumerate}
\columnbreak

\textbf{Три главные мысли}
\begin{itemize}[leftmargin=*,itemsep=2pt]
    \item Дрон — это автопилот: вы \emph{отдаёте команду} и \emph{ждёте события}, что она выполнена.
    \item В Lua ждут через \texttt{callback(event)}, в Python — опросом \texttt{while not \dots}.
    \item Координаты точек — метры в системе сцены, их показывает \key{Ctrl}+ЛКМ.
\end{itemize}
\end{multicols}

% ---------------------------------------------------------------- 1
\section{Быстрый старт}

\begin{center}
\begin{tikzpicture}
    \node[anchor=south west,inner sep=0] (img) at (0,0) {\shot{01-scene.png}};
    \begin{scope}[x={(img.south east)},y={(img.north west)}]
        \callout{0.075}{0.945}{1}
        \callout{0.905}{0.83}{2}
        \callout{0.595}{0.17}{3}
        \callout{0.12}{0.115}{4}
        \callout{0.343}{0.115}{5}
        \callout{0.86}{0.10}{6}
    \end{scope}
\end{tikzpicture}
\end{center}

\begin{multicols}{2}
\small
\begin{enumerate}[leftmargin=*,itemsep=1pt,label=\bubble{\arabic*}]
    \item Панели: код, телеметрия, рой, менеджер сцены, руководство, настройки.
    \item Компас осей (щелчок по оси — вид с неё) и окна камер дрона.
    \item \ui{Открыть редактор кода} — сюда пишется программа.
    \item Камера: над землёй, за дроном, от лица дрона (FPV), свободная.
    \item Расстановка объектов на сцене: дома, маркеры, машины.
    \item \ui{Запустить}, \ui{Стоп} и \ui{Сбросить} (дрон возвращается на старт).
\end{enumerate}
\end{multicols}

\textbf{Пять шагов до полёта:} открыть сайт $\to$ \ui{Открыть редактор кода} $\to$ выбрать язык в левом верхнем
списке (Lua, Python или Blockly) $\to$ вставить пример из раздела~2 $\to$ нажать \ui{Запустить}. Пока открыт
редактор, сцена скрыта: закройте его кнопкой \ui{Код} на панели слева — программа продолжит работать, и
вы увидите полёт.

\begin{tip}
Готовые сцены (жилой квартал, Геоскан Арена, гоночная трасса) — в \ui{Менеджере сцены}. Вращать камеру —
левой кнопкой мыши, сдвигать — правой, приближать — колёсиком.
\end{tip}

% ---------------------------------------------------------------- 2
\section{Первая программа на Lua}

Автопилот Пионера работает на событиях. Программа отдаёт команду и заканчивает, а когда дрон её выполнит,
автопилот сам вызывает вашу функцию \texttt{callback(event)} и сообщает, что произошло. Следующую команду
отдают \emph{там}, по нужному событию.

\begin{center}\shot[0.72\linewidth]{02-editor-lua.png}\\
\captionof{figure}{Программа в редакторе. Её текст — ниже, его можно скопировать из PDF.}\end{center}

\begin{lstlisting}[style=lua]
ap.push(Ev.MCE_PREFLIGHT)             -- запустить моторы

function callback(event)
    if event == Ev.ENGINES_STARTED then
        ap.push(Ev.MCE_TAKEOFF)       -- моторы запущены: взлетаем
    end
    if event == Ev.TAKEOFF_COMPLETE then
        ap.goToLocalPoint(2, 1, 1.5)  -- взлетели: летим в точку (x, y, z в метрах)
    end
    if event == Ev.POINT_REACHED then
        ap.push(Ev.MCE_LANDING)       -- долетели: садимся
    end
end
\end{lstlisting}

\begin{center}\shot[0.86\linewidth]{03-flight.png}\\
\captionof{figure}{Та же программа в полёте, камера \ui{Дрон}. Голубая линия — след полёта.}\end{center}

\begin{center}
\small
\begin{tabularx}{\linewidth}{@{}l X l X@{}}
\toprule
\multicolumn{2}{@{}l}{\textbf{Команды автопилоту} — \texttt{ap.push(\dots)}} & \multicolumn{2}{l@{}}{\textbf{События от автопилота} — в \texttt{callback}} \\
\midrule
\texttt{Ev.MCE\_PREFLIGHT} & запустить моторы & \texttt{Ev.ENGINES\_STARTED} & моторы запущены \\
\texttt{Ev.MCE\_TAKEOFF}   & взлететь         & \texttt{Ev.TAKEOFF\_COMPLETE} & взлёт закончен \\
\texttt{Ev.MCE\_LANDING}   & сесть            & \texttt{Ev.POINT\_REACHED}    & точка достигнута \\
\texttt{Ev.ENGINES\_DISARM}& выключить моторы & \texttt{Ev.COPTER\_LANDED}    & дрон сел \\
                           &                  & \texttt{Ev.LOW\_VOLTAGE2}     & батарея садится \\
\bottomrule
\end{tabularx}
\end{center}

\begin{warn}[Функция callback обязательна]
Даже если события не нужны, объявите пустую \texttt{function callback(event) end}. Без неё автопилот
выполнит только первую команду миссии, а остальные проигнорирует.
\end{warn}

% ---------------------------------------------------------------- 3
\section{Координаты: где точка}

\begin{center}\shot[0.78\linewidth]{04-coords.png}\\
\captionof{figure}{\key{Ctrl}+щелчок по земле или по объекту показывает координаты точки.}\end{center}

\begin{itemize}[itemsep=1pt]
    \item Координаты — в \textbf{метрах}, в системе сцены. Ось \textbf{Z смотрит вверх}, X и Y лежат на земле;
          их направления показывает компас в правом верхнем углу.
    \item Эти же числа понимают \texttt{ap.goToLocalPoint(x, y, z)} и \texttt{go\_to\_local\_point} в Python, и их же
          возвращают \texttt{Sensors.lpsPosition()} и \texttt{get\_local\_position\_lps()}: щёлкнули с \key{Ctrl} — подставили
          в программу.
    \item Показанное \textbf{Z — высота поверхности}, по которой щёлкнули. Лететь надо выше: над крышей
          дома на высоте 13~м задавайте, например, $z = 14$.
    \item Курс (\texttt{ap.updateYaw}, \texttt{yaw} в Python) — в \textbf{радианах}: $\pi/2 \approx 1.57$ — четверть оборота.
\end{itemize}

% ---------------------------------------------------------------- 4
\section{Несколько точек подряд}

Каждая точка заканчивается одним и тем же событием \texttt{POINT\_REACHED}. Чтобы понять, \emph{какая}
точка достигнута, заведите счётчик шагов — так сделаны официальные примеры Geoscan.

\begin{lstlisting}[style=lua]
local points = {{0, 0, 1}, {2, 0, 1}, {2, 2, 1.5}, {0, 2, 1}}
local step = 0

ap.push(Ev.MCE_PREFLIGHT)

local function nextPoint()
    step = step + 1
    if step <= #points then
        ap.goToLocalPoint(points[step][1], points[step][2], points[step][3])
    else
        ap.push(Ev.MCE_LANDING)            -- точки кончились — садимся
    end
end

function callback(event)
    if event == Ev.ENGINES_STARTED then ap.push(Ev.MCE_TAKEOFF) end
    if event == Ev.TAKEOFF_COMPLETE then nextPoint() end
    if event == Ev.POINT_REACHED then
        Timer.callLater(1, nextPoint)      -- секунду повисеть и дальше
    end
end
\end{lstlisting}

\begin{warn}[«Дрон не долетел до высоты»]
Частая ошибка: задать точку и в том же обработчике через \texttt{sleep(3)} отправить посадку. Дрон
поднимается примерно на 0.5~м/с, за 3~секунды он наберёт только полтора метра, и посадка отменит
подъём. Садиться нужно по \emph{следующему} \texttt{POINT\_REACHED}, когда точка действительно достигнута.
\end{warn}

\begin{tip}[Лететь медленнее или быстрее]
Четвёртый аргумент — время перелёта в секундах: \texttt{ap.goToLocalPoint(2, 1, 1.5, 8)} долетит за 8~секунд.
Без него дрон летит со скоростью из параметров автопилота.
\end{tip}

% ---------------------------------------------------------------- 5
\section{Таймеры, светодиоды, датчики}

\subsection{Таймеры}
\begin{itemize}[itemsep=1pt]
    \item \texttt{Timer.callLater(2, function() \dots\ end)} — выполнить один раз через 2~секунды.
    \item \texttt{Timer.new(0.5, function() \dots\ end):start()} — выполнять каждые 0.5~секунды;
          сохраните таймер в переменную, чтобы потом вызвать \texttt{:stop()}.
    \item \texttt{sleep(t)} тоже есть, но он останавливает всю программу — официальная документация
          советует таймеры.
\end{itemize}

\subsection{Светодиоды}
Цвет задаётся тремя числами от 0 до 1 (красный, зелёный, синий). На плате 4 светодиода (номера 0–3),
с модулем LED — 29.

\begin{lstlisting}[style=lua]
local leds = Ledbar.new(4)
local colors = {{1, 0, 0}, {0, 1, 0}, {0, 0, 1}}   -- красный, зелёный, синий
local i = 1

function callback(event) end

Timer.new(0.5, function()                          -- каждые полсекунды — следующий цвет
    for n = 0, 3 do leds:set(n, table.unpack(colors[i])) end
    i = i % #colors + 1
end):start()
\end{lstlisting}

\subsection{Датчики}
\begin{center}
\small
\begin{tabularx}{\linewidth}{@{}l X@{}}
\toprule
\texttt{local x, y, z = Sensors.lpsPosition()} & положение, м \\
\texttt{local vx, vy, vz = Sensors.lpsVelocity()} & скорость, м/с \\
\texttt{local d = Sensors.range()} & дальномер: расстояние до поверхности под дроном, м \\
\texttt{local h = Sensors.altitude()} & высота по барометру, м \\
\texttt{local yaw = Sensors.lpsYaw()} & курс, радианы \\
\texttt{local v = Sensors.battery()} & напряжение батареи, В \\
\texttt{local ch1, \dots, ch8 = Sensors.rc()} & каналы пульта, от $-1$ до $1$ \\
\bottomrule
\end{tabularx}
\end{center}

% ---------------------------------------------------------------- 6
\section{Python в браузере}

Выберите в списке языка \ui{Python}. Это тот же \texttt{pioneer\_sdk}, что и для настоящего Пионера:
программа, отлаженная здесь, подходит и дрону. Python ждёт не событий, а опросом — \texttt{while not \dots}.

\begin{center}\shot[0.66\linewidth]{05-editor-python.png}\end{center}

\begin{lstlisting}[style=py]
from pioneer_sdk import Pioneer

pioneer = Pioneer()

pioneer.arm()                  # запустить моторы
pioneer.takeoff()              # взлететь
pioneer.go_to_local_point(x=2, y=1, z=1.5, yaw=0)
while not pioneer.point_reached():
    pass                       # ждём, пока долетит
pioneer.land()                 # сесть

pioneer.close_connection()
\end{lstlisting}

\begin{itemize}[itemsep=1pt]
    \item \texttt{yaw} у \texttt{go\_to\_local\_point} обязателен (радианы); \texttt{yaw=0} — курс вдоль оси X.
    \item Точку можно отправить сразу после \texttt{takeoff()}: автопилот полетит к ней, когда закончит взлёт.
    \item \texttt{point\_reached()} возвращает \texttt{True} \emph{один раз} на каждую точку.
    \item Цвет в Python — от 0 до 255: \texttt{pioneer.led\_control(led\_id=255, r=0, g=255, b=0)} (255 — все светодиоды).
    \item Датчики: \texttt{get\_local\_position\_lps()}, \texttt{get\_dist\_sensor\_data()}, \texttt{get\_battery\_status()};
          передавайте \texttt{get\_last\_received=True}, чтобы повторное чтение не вернуло \texttt{None}.
\end{itemize}

% ---------------------------------------------------------------- 7
\section{Blockly: блоки вместо кода}

Выберите в списке языка \ui{Blockly} и собирайте программу из блоков под \ui{Начало программы}.
Блоки выполняются сверху вниз: \ui{Взлететь} и \ui{Лететь в точку} ждут, пока дрон закончит, — события и
счётчики за вас расставит сам Blockly.

\begin{center}\shot[0.78\linewidth]{06-blockly.png}\\
\captionof{figure}{Взлёт, три точки по оси X с зелёным светом и паузой в секунду, посадка.}\end{center}

\begin{itemize}[itemsep=1pt]
    \item Одна и та же программа из блоков превращается и в Lua, и в Python: переключите язык — и
          \ui{Скачать код} отдаст готовый текст, который можно дописать руками.
    \item Серые блоки недоступны в выбранном языке (например, ручная скорость есть только в Python,
          события автопилота — только в Lua). Подсказка появится, если навести курсор.
    \item \ui{Каждые N сек} — повтор по таймеру, \ui{Через N сек выполнить} — отложенное действие,
          \ui{Ждать, пока} — пауза до выполнения условия.
\end{itemize}

% ---------------------------------------------------------------- 8
\section{Камеры дрона}

\begin{center}\shot[0.86\linewidth]{07-cameras.png}\\
\captionof{figure}{Окна передней и нижней камер. Их кнопки — под компасом.}\end{center}

\begin{itemize}[itemsep=1pt]
    \item Окна включаются кнопками под компасом, перетаскиваются за заголовок; двойной щелчок
          по заголовку возвращает окно на место. \ui{На весь экран} в передней камере — вид от лица дрона на весь экран.
    \item Передняя камера — это кадр, который получает программа. В Lua снимок в загрузки браузера
          сохраняет \texttt{camera.requestMakeShot()}. В Python — как на настоящем дроне:
\end{itemize}
\begin{lstlisting}[style=py]
from pioneer_sdk import Camera
import cv2

camera = Camera()
frame = camera.get_cv_frame()          # кадр для OpenCV, цвета BGR
\end{lstlisting}

% ---------------------------------------------------------------- 9
\section{Python со своего компьютера}

Программу можно запускать не в браузере, а у себя, в обычном Python с официальным
\texttt{pioneer\_sdk} (\texttt{pip install pioneer-sdk}), — а летать будет дрон на сайте.

\begin{enumerate}[itemsep=1pt]
    \item Откройте симулятор, выберите дрон и включите у него \ui{Разрешить внешние команды}.
          Вкладку можно свернуть — полёт продолжится.
    \item В программе укажите адрес сайта вместо адреса дрона:
\end{enumerate}
\begin{lstlisting}[style=py]
from pioneer_sdk import Pioneer, Camera

pioneer = Pioneer(ip='simulator.sakoryagin.ru', mavlink_port=8001)

# Камера через интернет: один пакет от программы открывает домашний роутер для кадров.
class SimCamera(Camera):
    def connect(self):
        ok = super().connect()
        self.udp.sendto(b'1', (self.ip, self.port))
        return ok

camera = SimCamera(ip='simulator.sakoryagin.ru', port=18001)
\end{lstlisting}

% ---------------------------------------------------------------- 10
\section{Частые ошибки}

\begin{minipage}[t]{0.40\linewidth}
\vspace{0pt}
\shot{09-error.png}
\end{minipage}\hfill
\begin{minipage}[t]{0.56\linewidth}
\vspace{0pt}
Перед запуском симулятор проверяет программу и объясняет, что не так, в уведомлениях
(\ui{колокольчик} в заголовке редактора) — вместе с исправленным примером. Слева — что он скажет на
\texttt{ap.push(Ev.MCE\_PREFLIGHT)} и сразу \texttt{ap.push(Ev.MCE\_TAKEOFF)}.

\small
\begin{itemize}[leftmargin=*,itemsep=2pt]
    \item \textbf{Команды подряд без ожидания.} Взлёт — только после \texttt{ENGINES\_STARTED}, точка — после
          \texttt{TAKEOFF\_COMPLETE}, посадка — после \texttt{POINT\_REACHED}.
    \item \textbf{Посадка раньше точки.} \texttt{sleep} или таймер короче перелёта — дрон сядет, не долетев.
    \item \textbf{Нет \texttt{callback}.} Выполнится только первая команда миссии.
    \item \textbf{Все точки на одном событии.} Без счётчика шагов вторая точка затрёт первую
          (раздел~4).
    \item \textbf{Точка «на земле».} Z из \key{Ctrl}+ЛКМ — высота поверхности, летите выше неё.
    \item \textbf{Python: пропущен \texttt{yaw}.} На настоящем \texttt{pioneer\_sdk} это ошибка.
    \item \textbf{Скорость и логи.} Терминал (значок \texttt{>\_} в заголовке редактора) показывает
          \texttt{print} и сообщения автопилота.
\end{itemize}
\end{minipage}

% ---------------------------------------------------------------- 11
\section{Шпаргалка: Lua и Python}

\begin{center}
\small
\renewcommand{\arraystretch}{1.15}
\begin{tabularx}{\linewidth}{@{}>{\raggedright\arraybackslash}p{0.2\linewidth} >{\raggedright\arraybackslash}X >{\raggedright\arraybackslash}X@{}}
\toprule
\textbf{Действие} & \textbf{Lua} & \textbf{Python} \\
\midrule
Моторы          & \texttt{ap.push(Ev.MCE\_PREFLIGHT)} & \texttt{pioneer.arm()} \\
Взлёт           & \texttt{ap.push(Ev.MCE\_TAKEOFF)} & \texttt{pioneer.takeoff()} \\
Точка           & \texttt{ap.goToLocalPoint(x, y, z)} & \texttt{pioneer.go\_to\_local\_point(x=, y=, z=, yaw=)} \\
Ждать точку     & \texttt{if event == Ev.POINT\_REACHED} & \texttt{while not pioneer.point\_reached(): pass} \\
Курс            & \texttt{ap.updateYaw(рад)} & \texttt{yaw=} в \texttt{go\_to\_local\_point} \\
Посадка         & \texttt{ap.push(Ev.MCE\_LANDING)} & \texttt{pioneer.land()} \\
Пауза           & \texttt{Timer.callLater(t, f)} & \texttt{time.sleep(t)} \\
Светодиоды      & \texttt{leds:set(n, r, g, b)} — 0…1 & \texttt{pioneer.led\_control(led\_id, r, g, b)} — 0…255 \\
Позиция         & \texttt{Sensors.lpsPosition()} & \texttt{pioneer.get\_local\_position\_lps()} \\
Дальномер       & \texttt{Sensors.range()} & \texttt{pioneer.get\_dist\_sensor\_data()} \\
Батарея         & \texttt{Sensors.battery()} & \texttt{pioneer.get\_battery\_status()} \\
Снимок / кадр   & \texttt{camera.requestMakeShot()} & \texttt{Camera().get\_cv\_frame()} \\
Вывод           & \texttt{print(\dots)} & \texttt{print(\dots)} \\
\bottomrule
\end{tabularx}
\end{center}

\begin{tip}[Где искать остальное]
Кнопка со значком книги в заголовке редактора открывает справочник API с описанием и примером каждой
функции. Раздел \ui{Руководство} на панели слева — уроки с заданиями и автопроверкой.
\end{tip}

\end{document}
